Soient $n$ et $k$ des entiers strictement positifs tels que $k \geqslant n$ et $k - n$ est pair.

On suppose données $2n$ lampes numérotées de $1$ à $2n$ ; chacune peut être \emph{allumée} ou \emph{éteinte}.

Au début, toutes les lampes sont éteintes.

Une \emph{opération} consiste à allumer une lampe éteinte ou bien à éteindre une lampe allumée. On considère des
séquences constituées d'opérations successives.

Soit $N$ le nombre de séquences constituées de $k$ opérations et aboutissant à l'état où les lampes de $1$ à
$n$ sont allumées et les lampes de $n + 1$ à $2n$ sont éteintes.

Soit $M$ le nombre de séquences constituées de $k$ opérations et aboutissant à l'état où les lampes de $1$ à
$n$ sont allumées et les lampes de $n + 1$ à $2n$ sont éteintes, mais où les lampes de $n + 1$ à $2n$ n'ont
jamais été allumées.

Déterminer le rapport $N/M$.
